\chapter{Open-boundary matrix product states and sequential truncation}
\label{ch:one_dimensional_open_mps}

A vector on a chain is an open-boundary matrix product state of bond
dimension at most $D$ exactly when the Schmidt rank across every cut of
the chain is at most $D$. This chapter proves that characterization and
bounds the error made by truncating the Schmidt decompositions across
the cuts one after another. Chapter~\ref{ch:one_dimensional_chain} combines
this bound with the area law to prove Theorem~B.

\section{Cut ranks of matrix product states}

\begin{definition}[Open-boundary matrix product states and cut matrices]
    \label{def:chain_open_mps}
    \lean{
        TNLean.MPS.AreaLaw.IsOpenMPS,
        TNLean.MPS.AreaLaw.IsOpenMPS.smul,
        TNLean.MPS.AreaLaw.leftBlock,
        TNLean.MPS.AreaLaw.mem_leftBlock,
        TNLean.MPS.AreaLaw.leftBlock_eq_interval,
        TNLean.MPS.AreaLaw.cutMatrix,
        TNLean.MPS.AreaLaw.cutMatrix_apply
    }
    \leanok
    \uses{def:chain_model}
    Let $L,q,D$ be natural numbers. A vector $\psi$ on $L$ sites of local
    dimension $q$ is an \emph{open-boundary matrix product state of bond
        dimension at most $D$} when there are natural numbers $r_0,\ldots,r_L$
    with $r_0=r_L=1$ and $r_k\leq D$ for $0<k<L$, a scalar $c\in\mathbb C$,
    and matrices $A_i(s)\in\mathbb C^{r_i\times r_{i+1}}$ such that
    \begin{align}
        \psi(x_0,\ldots,x_{L-1})
         & =c\bigl(A_0(x_0)\cdots A_{L-1}(x_{L-1})\bigr)_{0,0}. \notag
    \end{align}
    The empty product is the identity on $\mathbb C$, so at $L=0$ this
    condition reads $\psi(\varnothing)=c$. Only the internal bond dimensions
    are bounded by $D$. Scaling $\psi$ scales $c$.
    The \emph{left block} of $k$ sites is the interval $[0,k)$, and the
    \emph{cut matrix} at $k$ has the configurations of the left block as rows
    and those of the remaining sites as columns; its rank is the Schmidt
    rank across the cut.
\end{definition}

\begin{lemma}[Products of one-site coefficients]
    \label{thm:chain_site_product_mps}
    \lean{TNLean.MPS.AreaLaw.isOpenMPS_prod,
        TNLean.MPS.AreaLaw.isOpenMPS_operator_entries_prod,
        TNLean.MPS.AreaLaw.isOpenMPS_operator_entries_rectKronecker,
        TNLean.MPS.AreaLaw.isOpenMPS_operator_entries_embed_one_site,
        TNLean.MPS.AreaLaw.isOpenMPS_operator_entries_of_supported_singleton}
    \leanok
    \uses{def:chain_open_mps,def:qc_local_circuit}
    Let $L,q$ be natural numbers and let $Q$ be an alphabet with
    $q$ elements. Given complex functions $f_i$ on $Q$, one at each site, the vector
    $\phi(w)=\prod_{i=0}^{L-1}f_i(w_i)$ has internal bond bound one.
    For a complex operator $M$, regard its entries as a vector on the
    paired local alphabet by putting
    \begin{align}
        v_M\bigl((x_0,y_0),\ldots,(x_{L-1},y_{L-1})\bigr)
         & =M_{x,y}. \notag
    \end{align}
    If $M=M_0\otimes\cdots\otimes M_{L-1}$ for one-site matrices
    $M_i$, then $v_M$ has internal bond bound one. The same is true
    for every operator supported on a single site, including the
    placement of an arbitrary one-site matrix with the identity at
    all other sites. Empty products and zero local dimension are
    included; a single-site placement presupposes a site in the chain.
    These identities give the bond-one ingredients for
    \cite[Proposition~5.7]{Arad2016RigorousRG}.
\end{lemma}
\begin{proof}\leanok
    \uses{def:chain_open_mps,lem:qc_placed_operators,
        thm:qc_supported_operator_placement}
    For a site product take the scalar multiplying the contraction
    to be one, all virtual dimensions to be one, and
    $A_i(s)_{0,0}=f_i(s)$. The unique virtual path gives $\phi(w)$.
    For tensor products of operators apply this construction to
    $f_i(x_i,y_i)=(M_i)_{x_i,y_i}$. An operator supported on one site
    is the placement of a one-site matrix; its other one-site factors
    are identities, so the same construction applies.
\end{proof}

\begin{lemma}[Bond-one states are scalar multiples of site products]
    \label{thm:chain_bond_one_products}
    \lean{TNLean.MPS.AreaLaw.IsOpenMPS.exists_eq_mul_prod,
        TNLean.MPS.AreaLaw.IsOpenMPS.exists_operator_entries_eq_mul_prod,
        TNLean.MPS.AreaLaw.isOpenMPS_one_iff_exists_eq_mul_prod}
    \leanok
    \uses{def:chain_open_mps}
    A complex vector $\psi$ on $L$ sites of local dimension $q$ has an
    open-boundary matrix product representation with internal bond
    bound one if and only if there are a scalar $c\in\mathbb C$ and
    one-site functions $f_i$ such that
    \begin{align}
        \psi(w) & =c\prod_{i=0}^{L-1}f_i(w_i). \notag
    \end{align}
    In particular, if the paired entries of a complex operator $M$
    have bond bound one, then there are one-site matrices $A_i$ and
    a scalar $c$ such that
    \begin{align}
        M_{x,y} & =c\prod_{i=0}^{L-1}(A_i)_{x_i,y_i}. \notag
    \end{align}
    The statements include the zero vector, empty alphabets and the
    empty chain. For the empty chain the product is one and the
    scalar $c$ is unrestricted.
    These identities give the bond-one ingredients for
    \cite[Proposition~5.7]{Arad2016RigorousRG}.
\end{lemma}
\begin{proof}\leanok
    \uses{def:chain_open_mps,thm:chain_site_product_mps}
    Choose an open-boundary representation. If one virtual bond has
    dimension zero, there is no virtual path and $\psi=0$; take
    $c=0$ and all $f_i=1$. Otherwise every virtual dimension is one.
    Put $f_i(s)=A_i(s)_{0,0}$. The unique virtual path gives
    $\psi(w)=c\prod_{i=0}^{L-1}f_i(w_i)$. The converse follows from
    Lemma~\ref{thm:chain_site_product_mps}, followed by multiplication
    by $c$. For the operator assertion, pair the two physical labels
    at each site and take $(A_i)_{a,b}=f_i(a,b)$. The pairing and
    unpairing are inverse bijections, so the row and column labels
    are recovered without conjugation.
\end{proof}

\begin{lemma}[Bounded internal cut ranks give a matrix product state]
    \label{thm:chain_mps_of_cut_rank}
    \lean{
        TNLean.MPS.AreaLaw.isOpenMPS_of_internal_cutMatrix_rank_le,
        TNLean.MPS.AreaLaw.isOpenMPS_of_cutMatrix_rank_le,
        Submodule.exists_fin_span_of_finrank_le,
        TNLean.MPS.AreaLaw.glueLeft,
        TNLean.MPS.AreaLaw.sliceMatrix,
        TNLean.MPS.AreaLaw.sliceMatrix_apply,
        TNLean.MPS.AreaLaw.rank_sliceMatrix_le,
        TNLean.MPS.AreaLaw.sliceSpace,
        TNLean.MPS.AreaLaw.finrank_sliceSpace_le,
        TNLean.MPS.AreaLaw.glueLeft_snoc,
        TNLean.MPS.AreaLaw.slice_mem_sliceSpace
    }
    \leanok
    Let $L,q,D$ be natural numbers and let $\psi$ be a vector on $L$ sites
    of local dimension $q$. If every internal cut matrix, at $0<k<L$, has
    rank at most $D$, then $\psi$ is an open-boundary matrix product state
    of bond dimension at most $D$. In particular, the conclusion holds
    when every cut matrix has rank at most $D$.
\end{lemma}
\begin{proof}
    \leanok
    \uses{def:chain_open_mps}
    Suppose first that $L\geq2$ and $D\geq1$. The two boundary cut matrices
    have rank at most $1\leq D$. Let $V_k$ be the column space of the cut
    matrix at $k$, and choose a spanning family $u_k$ of $D$ vectors.
    Slicing at a fixed value of site $k$ maps $V_{k+1}$ into $V_k$.
    Expressing these slices in the chosen spanning families gives square
    matrices $A_k(s)$ and boundary vectors $v_L,v_R$ with
    $\psi(x)=v_L^{\mathsf T}A_0(x_0)\cdots A_{L-1}(x_{L-1})v_R$.
    Absorb $v_L^{\mathsf T}$ into the first matrix and $v_R$ into the last;
    these are distinct matrices because $L\geq2$. Their rectangular
    dimensions have singleton endpoints and internal dimensions $D$.
    If $L\geq2$ and $D=0$, the cut matrix at $1$ vanishes, hence $\psi=0$.
    Set every internal bond dimension to zero and take zero site matrices
    and $c=0$. If $L=0$, take $c=\psi(\varnothing)$. If $L=1$, take
    $c=1$ and the $1\times1$ matrices $A_0(s)_{0,0}=\psi(s)$.
    The assertion about all cuts follows by restricting to internal cuts.
\end{proof}

\begin{theorem}[Cut ranks characterize open-boundary matrix product states]
    \label{thm:chain_cut_rank_of_mps}
    \lean{
        TNLean.MPS.AreaLaw.rank_cutMatrix_le_bondDim,
        TNLean.MPS.AreaLaw.IsOpenMPS.rank_cutMatrix_le,
        TNLean.MPS.AreaLaw.isOpenMPS_iff_internal_cutMatrix_rank_le
    }
    \leanok
    Let $\psi$ be a vector on $L$ sites of local dimension $q$.
    If $\psi$ is a scalar multiple of a rectangular open-boundary
    contraction with bond dimensions $d_0,\ldots,d_L$, then its cut
    matrix $M_k$ satisfies $\operatorname{rank}M_k\le d_k$ for
    $0\le k\le L$. In particular, $\psi$ is an open-boundary matrix
    product state of bond dimension at most $D$ if and only if
    $\operatorname{rank}M_k\le D$ at every internal cut $0<k<L$.
    No normalization is required. The endpoint ranks are at most one;
    when $L\le1$ there are no internal rank conditions, even for $D=0$.
    This is the single-cut rank observation used in the truncation-error
    lower bound of~\cite{Schuch2008EntropyScaling}; see also the sequential
    truncation proposition in~\cite[Part IIa]{TNLean2026OneDimensional}.
\end{theorem}
\begin{proof}
    \leanok
    \uses{def:chain_open_mps,thm:chain_mps_of_cut_rank}
    Split a virtual path at bond $k$. For each fixed
    bond index $\alpha\in\{0,\ldots,d_k-1\}$, the sums over the remaining
    virtual indices on the two sides are independent. Hence
    \begin{align}
        M_k(a,b)=c\sum_{\alpha=0}^{d_k-1}U_k(a,\alpha)V_k(\alpha,b).
    \end{align}
    Absorbing $c$ into $U_k$ factors $M_k$ through $\mathbb C^{d_k}$,
    so its rank is at most $d_k$. This remains valid when $d_k=0$:
    the sum is empty and the cut matrix vanishes. At the endpoints,
    $d_0=d_L=1$. The converse is
    Lemma~\ref{thm:chain_mps_of_cut_rank}, which includes empty and
    one-site chains.
\end{proof}

\begin{lemma}[Reduced states at a chain cut]
    \label{thm:chain_reduced_state_rank_error}
    \lean{TNLean.MPS.AreaLaw.IsOpenMPS.rank_regionState_leftBlock_le,
        TNLean.MPS.AreaLaw.rectangularTraceNorm_regionState_leftBlock_sub_le}
    \leanok
    For vectors $\psi,\phi$ on a finite chain, write
    $\rho_k=\operatorname{Tr}_{k+1,\ldots,L}|\psi\rangle\langle\psi|$ and
    $\sigma_k=\operatorname{Tr}_{k+1,\ldots,L}|\phi\rangle\langle\phi|$.
    Then $\|\rho_k-\sigma_k\|_1\le
        \bigl\||\psi\rangle\langle\psi|-|\phi\rangle\langle\phi|\bigr\|_1$.
    If $\phi$ is an open-boundary MPS of bond dimension at most $D$,
    then $\operatorname{rank}\sigma_k\le D$ at every internal cut.
    Neither assertion requires normalized vectors. These are the partial-trace
    and rank estimates used in the approximation lower bound
    of~\cite{Schuch2008EntropyScaling}.
\end{lemma}
\begin{proof}\leanok
    \uses{thm:chain_cut_rank_of_mps}
    In cut coordinates the reduced state is $CC^\dagger$, where $C$ is
    the coefficient matrix of $\phi$. Thus
    $\operatorname{rank}\sigma_k=\operatorname{rank}C\le D$.
    Reindexing configurations preserves the full trace norm, and partial
    trace contracts it. Apply these two facts to
    $|\psi\rangle\langle\psi|-|\phi\rangle\langle\phi|$.
\end{proof}

\begin{lemma}[Spectral-tail lower bound for MPS approximation]
    \label{thm:chain_spectral_tail_lower_bound}
    \lean{TNLean.MPS.AreaLaw.IsOpenMPS.sum_descEigenvalue_leftBlock_tail_le_traceNorm}
    \leanok
    Let $\psi$ be a unit vector on a chain of $L$ sites, and let $\phi$
    be an open-boundary matrix product state of bond dimension at most $D$.
    At an internal cut $0<k<L$, let $\lambda_1\ge\cdots\ge\lambda_M\ge0$
    be the eigenvalues of the prefix reduction $\rho_k$ of $\psi$. Then
    \begin{align}
        \sum_{j=D+1}^{M}\lambda_j
         & \le\bigl\||\psi\rangle\langle\psi|
        -|\phi\rangle\langle\phi|\bigr\|_1.
        \label{eq:chain_spectral_tail_lower_bound}
    \end{align}
    The sum is zero when $D\ge M$. Neither normalization of $\phi$ nor
    a small-error assumption is required. This is the approximation lower
    bound of~\cite[Eq.~(eps-lower), source lines 246--264]{Schuch2008EntropyScaling}.
\end{lemma}
\begin{proof}\leanok
    \uses{thm:chain_reduced_state_rank_error}
    Denote the trace error in (\ref{eq:chain_spectral_tail_lower_bound}) by $E$.
    The prefix reduction of $\phi$ has rank at most $D$.
    If $E<1$, the optimized pure-state error and Ky Fan overlap estimates
    bound the discarded mass by $E/2\le E$; see
    \cite[``Discarded marginal spectrum at small pure-state error'']{QICLean2026Blueprint}.
    If $E\ge1$, positivity and $\operatorname{tr}\rho_k=1$ bound the
    discarded mass by $1\le E$. This proves
    (\ref{eq:chain_spectral_tail_lower_bound}) in both cases.
\end{proof}

\begin{lemma}[Growth of chain cut ranks]
    \label{thm:chain_cut_rank_growth}
    \lean{
        TNLean.MPS.AreaLaw.rank_cutMatrix_add_le,
        TNLean.MPS.AreaLaw.rank_cutMatrix_le_add,
        TNLean.MPS.AreaLaw.rank_cutMatrix_le_pow_min_mul,
        TNLean.MPS.AreaLaw.rank_cutMatrix_le_pow_forty_mul
    }
    \leanok
    \uses{def:chain_open_mps}
    Let $L,q$ be natural numbers and let $\psi$ be a vector on $L$ sites
    of local dimension $q$. Write $r_k$ for the rank of its cut matrix at
    $k$. For natural numbers $a,d$ with $a+d\leq L$, one has
    \begin{align}
        r_{a+d} & \leq q^d r_a,
                & r_a           & \leq q^d r_{a+d}. \notag
    \end{align}
    Consequently, if $a\leq k\leq b\leq L$ and $r_a,r_b\leq R$, then
    \begin{align}
        r_k & \leq q^{\min\{k-a,b-k\}}R. \notag
    \end{align}
    If, in addition, $q\geq1$ and $b-a\leq80w$ for a natural number $w$,
    then
    \begin{align}
        r_k & \leq q^{40w}R. \notag
    \end{align}
    These bounds apply without a normalization or nonvanishing assumption
    on $\psi$; in particular, the first three bounds also allow $q=0$.
    They give the internal-cut estimate used when passing from blocks to
    sites in \cite[Part~II, proof of Theorem~B]{TNLean2026OneDimensional}.
\end{lemma}
\begin{proof}\leanok
    \uses{def:chain_open_mps}
    Regard the coefficients at two adjacent cuts as two flattenings of a
    tensor with factors corresponding to the left sites, one site, and
    the right sites. A rank factorization of the first flattening through
    a space of dimension $r$ gives a factorization of the second through
    a space of dimension $qr$: retain the one-site coordinate together
    with the factorization index. Thus moving one site to the left
    factor increases the rank by at most a factor $q$. Transposition gives
    the same estimate when moving it to the right factor. Iterating proves
    the first two inequalities, including zero-dimensional factors.
    Applying them from the nearer endpoint gives the minimum-exponent
    estimate. Finally,
    $\min\{k-a,b-k\}\leq40w$ because the two distances sum to $b-a\leq80w$.
    Monotonicity of $q^n$ for $q\geq1$ gives the last bound.
\end{proof}

\section{Sequential truncation}

\begin{lemma}[Sequential truncation in ordered orthonormal bases]
    \label{thm:chain_sequential_truncation_full}
    \lean{
        TNLean.MPS.AreaLaw.prefixBasisProjection,
        TNLean.MPS.AreaLaw.isOpenMPS_sequentialTruncation_prefixBasisProjection
    }
    \leanok
    \uses{def:chain_open_mps}
    Let $L,q,D$ be natural numbers and let $\psi$ be a vector on the chain
    of $L$ sites of local dimension $q$. For each internal cut $0<k<L$,
    let $N_k=\dim\mathcal H_{[0,k)}$ and choose an ordered orthonormal basis
    $b_{k,0},\ldots,b_{k,N_k-1}$ of $\mathcal H_{[0,k)}$. Let $\Pi_k$ be
    the orthogonal projection onto the span of its first $\min\{D,N_k\}$
    vectors, and put $P_k=\Pi_k\otimes1$. Define
    \begin{align}
        \Phi       & =P_1P_2\cdots P_{L-1}\psi, \notag                         \\
        \epsilon_k & =\operatorname{Re}\langle\psi,(1-P_k)\psi\rangle . \notag
    \end{align}
    Then $\Phi$ is an open-boundary matrix product state with every internal
    bond dimension at most $D$, and
    \begin{align}
        \|\Phi-\psi\|\le\sum_{k=1}^{L-1}\sqrt{\epsilon_k}. \notag
    \end{align}
    Empty products act as the identity, so this includes $L=0$ and $L=1$,
    without a choice of basis at a boundary cut. If $D\ge N_k$, then
    $\Pi_k=1$.

    This statement holds for every choice of ordered orthonormal bases.
    Proposition~7.1 of \cite[Part~IIa]{TNLean2026OneDimensional} is the
    instance in which $\psi$ is a unit vector and each basis is an
    eigenbasis of the reduced state at its cut, ordered by nonincreasing
    eigenvalue, with any choice of eigenvectors in a repeated-eigenvalue
    space. That instance, with $\epsilon_k$ identified as the sum of the
    discarded eigenvalues, is Lemma~\ref{thm:chain_sequential_truncation_eigen}.
\end{lemma}
\begin{proof}
    \leanok
    \uses{def:chain_open_mps, thm:chain_mps_of_cut_rank,
        lem:chain_sequential_truncation}
    Each $P_k$ is an orthogonal projection, hence a contraction. Telescoping
    gives
    $\Phi-\psi=\sum_k P_1\cdots P_{k-1}(P_k-1)\psi$,
    and $\|(1-P_k)\psi\|^2=\epsilon_k$, proving the error bound.
    At cut $k$, the vector $P_k\cdots P_{L-1}\psi$ lies in
    $\operatorname{ran}\Pi_k\otimes\mathcal H_{[k,L)}$, so its cut rank
    is at most $\min\{D,N_k\}\le D$. The remaining factors act on the
    left of that cut and cannot increase this rank.
    Lemma~\ref{thm:chain_mps_of_cut_rank} gives the stated representation
    of $\Phi$.
\end{proof}

\begin{lemma}[Sequential truncation onto leading eigenvectors]
    \label{thm:chain_sequential_truncation_eigen}
    \lean{TNLean.MPS.AreaLaw.isOpenMPS_sequentialTruncation_leadingEigenprojection}
    \leanok
    \uses{def:chain_open_mps}
    Let $L,q,D$ be natural numbers and let $\psi$ be a unit vector on the
    chain of $L$ sites of local dimension $q$. For each internal cut
    $0<k<L$, let $\rho_k$ be the reduced state of $\psi$ on
    $\mathcal H_{[0,k)}$, let $\Pi_k$ be the projection onto the span of
    the eigenvectors of $\rho_k$ belonging to its $\min\{D,N_k\}$ largest
    eigenvalues, counted with multiplicity, and let $\epsilon_k(D)$ be the
    sum of the remaining eigenvalues. With $P_k=\Pi_k\otimes1$, the vector
    $\Phi=P_1P_2\cdots P_{L-1}\psi$ is an open-boundary matrix product state
    with every internal bond dimension at most $D$, and
    \begin{align}
        \|\Phi-\psi\|\le\sum_{k=1}^{L-1}\sqrt{\epsilon_k(D)}. \notag
    \end{align}
    This is Proposition~7.1 of \cite[Part~IIa]{TNLean2026OneDimensional}.
    Within a repeated eigenvalue the eigenvectors are taken from one fixed
    orthonormal eigenbasis; every other choice is covered by
    Lemma~\ref{thm:chain_sequential_truncation_full}.
\end{lemma}
\begin{proof}
    \leanok
    \uses{def:chain_open_mps, thm:chain_mps_of_cut_rank,
        lem:chain_sequential_truncation}
    Each $\Pi_k$ is an orthogonal projection of rank
    $\min\{D,N_k\}\le D$, and the trace of $(1-\Pi_k)\rho_k$ is
    $\operatorname{tr}\rho_k$ minus the sum of the $\min\{D,N_k\}$ largest
    eigenvalues, that is $\epsilon_k(D)$, since
    $\operatorname{tr}\rho_k=\|\psi\|^2=1$. The estimates of
    Lemma~\ref{lem:chain_sequential_truncation} hold for any orthogonal
    projections on the first $k$ sites: the cut rank of $\Phi$ at cut $k$ is
    at most the rank of $\Pi_k$, and
    $\|\Phi-\psi\|\le\sum_k\sqrt{\operatorname{tr}((1-\Pi_k)\rho_k)}$,
    because $\|(1-P_k)\psi\|^2=\operatorname{tr}((1-\Pi_k)\rho_k)$.
    Lemma~\ref{thm:chain_mps_of_cut_rank} turns the rank bound into the
    stated representation, and substituting $\epsilon_k(D)$ for each trace
    gives the error bound. The argument is the proof of
    Lemma~\ref{thm:chain_sequential_truncation_full}, applied to the
    eigenprojections directly.
\end{proof}

\begin{theorem}[Squared error of sequential truncation]
    \label{thm:chain_sequential_truncation_sq}
    \lean{
        TNLean.MPS.AreaLaw.norm_sequentialTruncation_leadingEigenprojection_sub_sq_le,
        TNLean.MPS.AreaLaw.norm_sub_sequentialTruncation_sq_le}
    \leanok
    \uses{thm:chain_sequential_truncation_eigen}
    In the setting of Lemma~\ref{thm:chain_sequential_truncation_eigen},
    \begin{align}
        \|\Phi-\psi\|^2\le\sum_{k=1}^{L-1}\epsilon_k(D). \notag
    \end{align}
    Verstraete and Cirac~\cite{VerstraeteCirac2006Faithfully} state this
    bound with a factor $2$ on the right; the form here implies theirs.
\end{theorem}
\begin{proof}
    \leanok
    Write $\Phi_0=\psi$ and $\Phi_j=P_{L-j}\Phi_{j-1}$, so that
    $\Phi=\Phi_{L-1}$. For an orthogonal projection $P$,
    \begin{align}
        \psi-P\Phi_{j-1} & =P(\psi-\Phi_{j-1})+(1-P)\psi, \notag
    \end{align}
    and the two terms are orthogonal, since $P^*=P=P^2$ gives
    \begin{align}
        \langle P(\psi-\Phi_{j-1}),(1-P)\psi\rangle
         & =\langle\psi-\Phi_{j-1},(P-P^2)\psi\rangle=0. \notag
    \end{align}
    Hence
    $\|\psi-\Phi_j\|^2\le\|\psi-\Phi_{j-1}\|^2+\|(1-P_{L-j})\psi\|^2$,
    and $\|(1-P_k)\psi\|^2=\operatorname{tr}((1-\Pi_k)\rho_k)=\epsilon_k(D)$.
    Summing over the steps gives the bound.
\end{proof}

\begin{lemma}[Sequential truncation at an entropy threshold]
    \label{lem:chain_sequential_truncation}
    \lean{
        Matrix.schmidtCoeffMatrix_kronecker_one_mulVec,
        Matrix.toEuclideanLin_one_apply,
        Matrix.toEuclideanLin_mul_apply,
        Entropy.localLift_sub,
        Entropy.exists_localLift_eq_localLift_of_subset,
        Entropy.norm_sq_toEuclideanLin_localLift_eq_trace,
        Entropy.norm_toEuclideanLin_localLift_le,
        Entropy.trace_regionState_eq_one,
        Matrix.PosSemidef.trace_complement_spectralProjectionGE_le_surprisalTail,
        Matrix.PosSemidef.rank_spectralProjectionGE_exp_le,
        TNLean.MPS.AreaLaw.cutMatrix_localLift,
        TNLean.MPS.AreaLaw.rank_cutMatrix_localLift_le,
        TNLean.MPS.AreaLaw.rank_cutMatrix_zero_le,
        TNLean.MPS.AreaLaw.rank_cutMatrix_le_one_of_le,
        TNLean.MPS.AreaLaw.LeftBlockMatrix,
        TNLean.MPS.AreaLaw.cutTruncation,
        TNLean.MPS.AreaLaw.sequentialTruncation,
        TNLean.MPS.AreaLaw.leftBlock_mono,
        TNLean.MPS.AreaLaw.norm_sub_sequentialTruncation_le,
        TNLean.MPS.AreaLaw.exists_sequentialTruncation_eq_localLift,
        TNLean.MPS.AreaLaw.rank_cutMatrix_sequentialTruncation_le,
        TNLean.MPS.AreaLaw.exists_sequentialTruncation
    }
    \leanok
    \uses{def:chain_open_mps}
    Let $\psi$ be a unit vector on $L$ sites and let $w$ be a real number. For
    each cut $k$,
    $1\le k\le L-1$, let $S_k$ be the entropy of the first $k$ sites, let
    $\Pi_k$ be the spectral projector of their reduced state onto the
    eigenvalues at least $e^{-(S_k+w)}$, an operator on the first $k$ sites,
    and let $P_k=\Pi_k\otimes1$ be its extension by the identity on the
    remaining sites. Then the regional projector $\Pi_k$ has rank at most
    $e^{S_k+w}$, the mass $\|(1-P_k)\psi\|^2$ discarded by $P_k$ is at most
    the surprisal tail of the reduced state beyond $w$, and the vector
    $\Phi=P_1P_2\cdots P_{L-1}\psi$ has cut rank at most
    $\max\{1,e^{S_k+w}\}$ across every cut $k$ of the chain, the trivial
    cuts at its ends included; it satisfies
    \begin{align}
        \|\psi-\Phi\|\le\sum_{k=1}^{L-1}\sqrt{\text{tail}_k(w)} .
    \end{align}
    This is the entropy-threshold form used in the proof of Theorem~B
    of \cite{TNLean2026OneDimensional}. The prescribed-dimension statement,
    with the actual discarded masses, is stated separately in
    Lemma~\ref{thm:chain_sequential_truncation_full}.
    A common entropy slack and a prescribed number of retained eigenvectors
    need not give the same projectors; see
    \cite{gap:tnlean_one_dimensional_threshold_truncation}.
\end{lemma}
\begin{proof}
    \leanok
    \uses{def:chain_open_mps}
    Write $M_k(\psi)$ for the cut matrix of $\psi$ at $k$, with the
    configurations of the first $k$ sites as rows. For an operator $\Pi$ on
    the first $k$ sites, the extension $\Pi\otimes1$ acts on it by left
    multiplication,
    \begin{align}
        M_k\bigl((\Pi\otimes1)\psi\bigr)=\Pi\,M_k(\psi),
    \end{align}
    so the cut rank of $(\Pi\otimes1)\psi$ is at most the rank of $\Pi$; the
    rank of the extended operator itself, which is the rank of $\Pi$ times the
    dimension of the remaining sites, plays no role. The projector on the
    smallest block is
    applied last, and the product of the projectors on blocks of size at most
    $k$ is one extension of an operator on the first $k$ sites ending in
    $\Pi_k$. The distance follows by induction from
    $\psi-P_k\Phi'=(1-P_k)\psi+P_k(\psi-\Phi')$ and the contraction property
    of a projector, and $\|(1-P_k)\psi\|^2$ is the trace of $1-\Pi_k$ against
    the reduced state, the mass of the eigenvalues below the threshold, which
    is at most the surprisal tail.
\end{proof}
